\chapter{Głębokie uczenie w~widzeniu komputerowym}

Rozdział przedstawia teoretyczne podstawy głębokich sieci neuronowych wykorzystywanych w~zadaniach widzenia komputerowego (ang.\ \emph{Computer Vision}). W~rozdziale przedstawiono rozwój architektur wykorzystywanych w~widzeniu komputerowym – od klasycznych sieci konwolucyjnych do modeli opartych na mechanizmie attention. Następnie opisano wybrane modele detekcji obiektów, w~szczególności architektury YOLO oraz RT-DETR, które zostały wykorzystane w~części eksperymentalnej pracy. Na końcu rozdziału przedstawiono definicje standardowych metryk wykorzystywanych do ewaluacji jakości i~zdolności generalizacji modeli wizyjnych.

\section{Podstawy sztucznych sieci neuronowych}

Sztuczne sieci neuronowe (ang.\ \emph{Artificial Neural Networks}) to modele matematyczne inspirowane budową i~funkcjonowaniem biologicznego układu nerwowego. Sieci neuronowe są obecnie jedną z~podstawowych metod stosowanych w~głębokim uczeniu. Jak wskazują Goodfellow i~in.\ w~podręczniku akademickim \cite{Goodfellow2016}, ich głównym zadaniem jest aproksymacja złożonych, nieliniowych funkcji, co pozwala na mapowanie danych wejściowych (np.\ pikseli obrazu) na pożądane wyjście (np.\ etykietę klasy lub współrzędne obiektu).

\subsection{Architektura sztucznego neuronu i~perceptron wielowarstwowy}

Podstawowym elementem obliczeniowym każdej sieci jest sztuczny neuron. Przyjmuje on wektor danych wejściowych $x = [x_1, x_2, \dots, x_n]$, gdzie każda z~wartości mnożona jest przez przypisaną jej wagę synaptyczną $w_i$. Suma ważona wejść powiększana jest o~wyraz wolny, zwany obciążeniem (ang.\ \emph{bias}, $b$). Ostateczny wynik przekazywany jest przez nieliniową funkcję aktywacji $f$, co można zapisać równaniem:
\begin{equation}
y = f\left(\sum_{i=1}^{n} w_i x_i + b\right)
\end{equation}

Funkcja aktywacji jest kluczowym elementem sieci, ponieważ wprowadza nieliniowość do systemu, umożliwiając rozwiązywanie problemów separowalnych nieliniowo. Obecnie standardem w~głębokich sieciach neuronowych jest funkcja ReLU (ang.\ \emph{Rectified Linear Unit}), zdefiniowana jako $f(x) = \max(0, x)$, która – co udowodnili w~swoich badaniach Nair i~Hinton \cite{Nair2010} – upraszcza obliczenia oraz w~znacznym stopniu ogranicza problem zanikania gradientu podczas uczenia głębokich sieci.

Neurony łączone są w~struktury zwane warstwami. Najprostszą topologią jest perceptron wielowarstwowy (ang.\ \emph{Multilayer Perceptron}), w~którym występują: warstwa wejściowa, jedna lub więcej warstw ukrytych (ang.\ \emph{hidden layers}) oraz warstwa wyjściowa. W~klasycznym wariancie MLP każdy neuron z~danej warstwy połączony jest z~każdym neuronem warstwy następnej, tworząc warstwę gęstą (ang.\ \emph{fully-connected layer}).

\subsection{Mechanizmy uczenia sieci neuronowych}

Uczenie sieci neuronowej to proces optymalizacyjny polegający na iteracyjnym dostrajaniu jej parametrów (wag i~obciążeń) w~taki sposób, aby zminimalizować błąd popełniany przez model. Różnica między predykcją sieci a~rzeczywistą (oczekiwaną) wartością zdefiniowana jest za pomocą funkcji straty (ang.\ \emph{Loss function}).

Kluczowym algorytmem umożliwiającym trenowanie wielowarstwowych sieci neuronowych jest propagacja wsteczna błędu (ang.\ \emph{Backpropagation}). Jak opisali to Rumelhart i~in.\ w~swojej przełomowej pracy \cite{Rumelhart1986}, metoda ta wykorzystuje regułę łańcuchową rachunku różniczkowego do obliczenia gradientu funkcji straty względem każdego parametru w~sieci. Informacja o~błędzie propagowana jest od warstwy wyjściowej, poprzez warstwy ukryte, aż do wejścia.

Mając wyznaczone gradienty, parametry są aktualizowane za pomocą algorytmów optymalizacyjnych z~rodziny spadku wzdłuż gradientu (ang.\ \emph{Gradient Descent}). W~nowoczesnych zastosowaniach dominuje optymalizator Adam \cite{Kingma2014} (ang.\ \emph{Adaptive Moment Estimation}), który dynamicznie dostosowuje tempo uczenia (ang.\ \emph{learning rate}) dla każdej wagi z~osobna, uwzględniając momenty pierwszego i~drugiego rzędu gradientów, co pozwala na stabilniejszą i~szybszą zbieżność modelu.

\subsection{Konwolucyjne sieci neuronowe i~operacja splotu}

Klasyczne, gęste sieci MLP są nieefektywne w~przetwarzaniu obrazów. Spłaszczenie dwuwymiarowego obrazu do jednowymiarowego wektora niszczy relacje przestrzenne między pikselami. Dodatkowo, jak zauważa Goodfellow \cite{Goodfellow2016}, pełne połączenia przy obrazach o~wysokiej rozdzielczości powodują znaczny wzrost liczby parametrów, co istotnie zwiększa ryzyko zjawiska przeuczenia (ang. \emph{overfitting}).

Rozwiązaniem tego problemu, zaprezentowanym m.in.\ przez Yanna LeCuna \cite{LeCun1998}, stały się Konwolucyjne Sieci Neuronowe (ang.\ \emph{Convolutional Neural Networks}, CNN). Wykorzystują one operację matematyczną splotu (ang.\ \emph{convolution}), w~której niewielka macierz wag, nazywana filtrem lub jądrem (ang.\ \emph{kernel}), przesuwa się po powierzchni obrazu wejściowego. Filtr ten wykonuje mnożenie punktowe z~lokalnym obszarem obrazu, tworząc mapę cech (ang.\ \emph{feature map}).

Zastosowanie operacji splotu pozwala na współdzielenie parametrów oraz ograniczenie liczby połączeń pomiędzy neuronami. W~typowej architekturze CNN warstwy splotowe przeplatane są nieliniowymi funkcjami aktywacji oraz warstwami łączącymi (ang.\ \emph{pooling layers}, np.\ \emph{Max Pooling}). Operacja ta służy do redukcji wymiarowości przestrzennej obrazu (ang.\ \emph{downsampling}), co zmniejsza koszt obliczeniowy oraz zapewnia sieci pewien stopień niezmienniczości na przesunięcia obiektów w~kadrze. W~głębokich warstwach sieci CNN mapy cech stają się mniejsze przestrzennie, ale rośnie ich liczba (liczba kanałów), co pozwala modelowi na detekcję złożonych, abstrakcyjnych struktur semantycznych.


\section{Wybrane architektury detekcji obiektów}

Weryfikacja metod poprawy jakości widzenia komputerowego na małym zbiorze danych wymaga zastosowania współczesnych architektur. W~ramach niniejszej pracy do analizy wytypowane zostały dwie klasy modeli reprezentujące odmienne podejścia do detekcji obiektów: detektory jednoetapowe z~rodziny YOLO oraz architektura RT-DETR oparta na wizyjnych mechanizmach uwagi (transformerach). Obie architektury wykorzystywane są do zadania klasycznej detekcji obiektów, polegającego na precyzyjnej lokalizacji obiektu na obrazie poprzez wyznaczenie prostokątnej ramki otaczającej (ang.\ \emph{bounding box}) oraz przypisaniu mu odpowiedniej klasy.

\subsection{Architektury z~rodziny YOLO}
Architektury z~rodziny YOLO (ang.\ \emph{You Only Look Once}) należą do grupy detektorów jednoetapowych (ang.\ \emph{one-stage detectors}). W~przeciwieństwie do modeli dwuetapowych (takich jak Faster R-CNN, wykorzystujących dodatkową sieć propozycji regionów), YOLO traktuje detekcję obiektów jako pojedynczy problem regresji. Pierwotna koncepcja, zaproponowana przez Redmona i~in.\ \cite{Redmon2016}, zakładała podział obrazu na siatkę $S \times S$, gdzie komórka zawierająca środek obiektu jest bezpośrednio odpowiedzialna za predykcję jego współrzędnych oraz klasy. Takie podejście, opierające się na globalnym wnioskowaniu o~obrazie, pozwoliło na znaczne zwiększenie szybkości działania modeli detekcji obiektów.

Nowsze wersje modeli YOLO zostały mocno rozbudowane względem wcześniejszych architektur i~obecnie składają się z~trzech głównych części: szkieletu (ang. \emph{backbone}) odpowiedzialnego za ekstrakcję cech, szyi (ang. \emph{neck}) służącej do ich agregacji oraz głowy (ang. \emph{head}) realizującej końcową detekcję obiektów. W~części ekstrakcyjnej często stosowane są struktury CSPNet, które, jak wykazali Wang i~in.\ \cite{Wang2020}, ograniczają koszt obliczeniowy modelu i~pomagają w~stabilniejszym uczeniu sieci. Z~kolei w~module szyi wykorzystuje się warianty architektury PANet \cite{Liu2018}, pozwalające na efektywne łączenie informacji przestrzennych z~różnych poziomów sieci.

Ze względu na wykorzystanie w~badaniach opisanych w~kolejnych rozdziałach pracy, przedstawiono architektury dwóch generacji z~tej rodziny: zoptymalizowanej wersji \textbf{YOLOv8} \cite{Jocher2023_YOLOv8} oraz nowszej, eksperymentalnej generacji \textbf{YOLOv26} \cite{Jocher2024_YOLOv26}. Model YOLOv8, wprowadzony przez zespół Ultralytics, spopularyzował w~rozwiązaniach przemysłowych podejście detektorów typu \emph{anchor-free} (wolnych od zdefiniowanych z~góry kotwic). Wprowadził on ponadto zaawansowane funkcje straty dla regresji położenia (takie jak \emph{Distribution Focal Loss}, DFL \cite{Li2020}), co znacząco poprawiło precyzję wykrywania małych obiektów.

Obie omawiane generacje czerpią z~technik optymalizacyjnych określanych zbiorczo jako \emph{Bag of Freebies} \cite{Wang2023}. Obejmują one m.in.\ zoptymalizowane funkcje straty typu CIoU (ang.\ \emph{Complete Intersection over Union}). Funkcja CIoU uwzględnia nie tylko nakładanie się ramek, ale również odległość między ich środkami oraz proporcje boków \cite{Zheng2020}. Mimo wspólnych założeń projektowych, architektura YOLOv26 wprowadza szereg zmian względem modelu YOLOv8, dotyczących przede wszystkim organizacji bloków obliczeniowych oraz sposobu przetwarzania cech.

Jedną z~głównych różnic pomiędzy YOLOv8 a~YOLOv26 jest sposób organizacji bloków odpowiedzialnych za ekstrakcję i~propagację cech. W~modelu YOLOv8 głównym elementem budulcowym odpowiedzialnym za ekstrakcję był moduł C2f (ang. \emph{Cross Stage Partial Bottleneck with 2 convolutions}). Choć rozwiązanie to sprzyjało stabilnemu procesowi uczenia, charakteryzował się znaczną pojemnością informacyjną i~dużą liczbą parametrów. W~architekturze YOLOv26 moduły typu C2f zostały zastąpione przez znacznie nowocześniejsze, zoptymalizowane bloki konwolucyjne (np. czerpiące z~koncepcji struktury C3k2 lub C2f-Star). Zmiana ta pozwoliła wyeliminować tzw. "wąskie gardła" obliczeniowe oraz zredukować niepotrzebne powielanie informacji o~gradiencie w~głębokich warstwach sieci, co przyczyniło się do zmniejszenia kosztu obliczeniowego modelu.

Kolejną istotną różnicą jest sposób wykorzystania mechanizmów uwagi (ang. \emph{Attention}). O~ile w~YOLOv8 mechanizmy te ograniczały się zazwyczaj do prostej agregacji przestrzennej na końcu szkieletu sieci (np. poprzez moduły SPPF), YOLOv26 wykorzystuje bardziej rozbudowane mechanizmy uwagi bezpośrednio w~strukturę ekstrakcji. Poprzez zastosowanie lżejszych mechanizmów (przypominających koncepcje Partial Self-Attention czy CBAM), model może skuteczniej akcentować cechy istotne z~punktu widzenia procesu detekcji już na wczesnym etapie jego przetwarzania, bez ponoszenia bardzo wysokich kosztów obliczeniowych typowych dla klasycznych transformatorów wizyjnych.

Z~perspektywy niniejszej pracy istotne jest, że modele YOLOv8 i~YOLOv26 reprezentują odmienne podejścia do projektowania współczesnych detektorów obiektów. W~architekturze YOLOv26 większy nacisk położono na efektywność obliczeniową oraz optymalizację przepływu informacji w~sieci. Porównanie obu modeli pozwala przeanalizować wpływ tych różnic architektonicznych na skuteczność detekcji obiektów w~warunkach ograniczonej dostępności danych treningowych.


\subsection{Architektura RT-DETR}

Alternatywnym podejściem do detekcji obiektów jest wykorzystanie transformerów wizyjnych, zapoczątkowane przez architekturę DETR (ang.\ \emph{DEtection TRansformer}) \cite{Carion2020}. Modele te całkowicie eliminują konieczność stosowania heurystycznych, trudnych w~dostrajaniu komponentów (takich jak generatory kotwic czy algorytmy tłumienia wartości niemaksymalnych -- NMS). Zamiast tego proces detekcji traktowany jest jako operacja bezpośredniego przewidywania zbioru, oparta na globalnym dopasowywaniu predykcji do etykiet rzeczywistych przy użyciu algorytmu dopasowania dwudzielnego (ang.\ \emph{bipartite matching}).

Najnowsze badania wskazują, że architektury wywodzące się z~rodziny DETR posiadają duży potencjał w~zadaniach o~ograniczonej dostępności danych (ang.\ \emph{few-shot object detection}) oraz podczas trenowania z~wykorzystaniem niedoskonałych pseudoetykiet, co wykazali Nguyen i~in.\ w~modelu Counting-DETR \cite{Nguyen2022}.

Głównym ograniczeniem pierwotnego DETR był jednak wysoki koszt obliczeniowy mechanizmu uwagi dla obrazów o~wysokiej rozdzielczości oraz wolna konwergencja podczas procesu uczenia. Odpowiedzią na te wyzwania jest architektura RT-DETR (ang.\ \emph{Real-Time DEtection TRansformer}) \cite{Lv2023}. W~modelu tym zastosowano hybrydowy koder łączący operacje konwolucyjne z~mechanizmem attention, co pozwala ograniczyć koszt obliczeniowy przy zachowaniu zdolności modelowania globalnych zależności.

Konstrukcja RT-DETR rozdziela przetwarzanie cech w~obrębie pojedynczej skali od ich integracji pomiędzy różnymi skalami. Operacje lokalne realizowane są przez bloki konwolucyjne, natomiast za fuzję informacji odpowiada mechanizm attention. Dodatkowo architektura wykorzystuje mechanizm wyboru zapytań uwzględniający przewidywaną jakość lokalizacji obiektów (ang.\ \emph{IoU-aware query selection}), co poprawia jakość reprezentacji przekazywanych do dekodera transformerowego.

Dzięki tym rozwiązaniom RT-DETR osiąga konkurencyjną dokładność przy jednoczesnym zachowaniu możliwości pracy w~czasie rzeczywistym, stanowiąc interesującą alternatywę dla współczesnych architektur detekcyjnych opartych zarówno na klasycznych sieciach konwolucyjnych, jak i~innych wariantach transformerów wizyjnych \cite{Lv2023}.

\section{Podsumowanie rozdziału}

Zaprezentowane w~niniejszym rozdziale architektury głębokich sieci neuronowych, w~szczególności modele z~rodziny YOLO oraz oparte na transformerach detektory RT-DETR, należą obecnie do najważniejszych rozwiązań wykorzystywanych w~zadaniach detekcji obiektów. Pomimo wysokiej skuteczności ich działanie jest silnie uzależnione od jakości i~liczebności dostępnych danych treningowych. W~warunkach ograniczonej dostępności danych modele te mogą wykazywać zwiększoną podatność na przeuczenie, co prowadzi do pogorszenia zdolności generalizacji na nowe przykłady. Przedstawione podstawy teoretyczne stanowią punkt wyjścia do dalszej analizy metod mających na celu poprawę skuteczności modeli detekcyjnych trenowanych na małych zbiorach danych. 
